\subsection{Measurable Spaces}
Measure theory is a branch of mathematics that studies measures and integrals on general spaces. A measure can be thought of as a generalisation of a length, area, or volume function to an arbitrary one-or multi-dimensional space.

In what follows, we let \(E\) be an arbitrary set, e.g., the interval \((0,1]\) or the square \((0,1] \times (0,1]\). We wish to assign a "measure" (think of "length" or "area") to certain subsets of \(E\). We can apply the usual set operations to subsets of \(E\), as illustrated in \Cref{fig:VennDiagrams}.

\begin{figure}[H]
	\centering
	\begin{subfigure}{0.24\textwidth}\centering
		\begin{tikzpicture}[thick,set/.style={circle,minimum size=1.5cm,draw=black},shaded/.style={fill=airforceblue!60}]
			\draw (-1.2,-1.2) rectangle (2.4,1.2);
			\begin{scope}
				\clip (0,0) circle (0.75cm);
				\fill[shaded] (1.2,0) circle (0.75cm);
			\end{scope}
			\node[set] at (0,0) {$A$};
			\node[set] at (1.2,0) {$B$};
		\end{tikzpicture}
		\caption{$A \cap B$}
	\end{subfigure}\hfill
	\begin{subfigure}{0.24\textwidth}\centering
		\begin{tikzpicture}[thick,set/.style={circle,minimum size=1.5cm,draw=black},shaded/.style={fill=airforceblue!60}]
			\draw (-1.2,-1.2) rectangle (2.4,1.2);
			\fill[shaded] (0,0) circle (0.75cm);
			\fill[shaded] (1.2,0) circle (0.75cm);
			\node[set] at (0,0) {$A$};
			\node[set] at (1.2,0) {$B$};
		\end{tikzpicture}
		\caption{$A \cup B$}
	\end{subfigure}\hfill
	\begin{subfigure}{0.24\textwidth}\centering
		\begin{tikzpicture}[thick,set/.style={circle,minimum size=1.5cm,draw=black},shaded/.style={fill=airforceblue!60}]
			\fill[shaded] (-1.2,-1.2) rectangle (2.4,1.2);
			\fill[white] (0.6,0) circle (0.75cm);
			\draw (-1.2,-1.2) rectangle (2.4,1.2);
			\node[set] at (0.6,0) {$A$};
		\end{tikzpicture}
		\caption{$A^c$}
	\end{subfigure}\hfill
	\begin{subfigure}{0.24\textwidth}\centering
		\begin{tikzpicture}[thick,set/.style={circle,minimum size=1.5cm,draw=black},shaded/.style={fill=airforceblue!60}]
			\draw (-1.2,-1.2) rectangle (2.4,1.2);
			\begin{scope}
				\clip (0,0) circle (0.75cm);
				\fill[shaded] (0,0) circle (0.75cm);
				\fill[white] (1.2,0) circle (0.75cm);
			\end{scope}
			\node[set] at (0,0) {$A$};
			\node[set] at (1.2,0) {$B$};
		\end{tikzpicture}
		\caption{$A \setminus B$}
	\end{subfigure}
	\caption{Venn diagrams of set operations.}
	\label{fig:VennDiagrams}
\end{figure}

\begin{defn}{Algebra}{AlgebraOnE}
	A collection $\mathcal{E}_0$ of subsets of $E$ is called an \textit{algebra on $E$} (or algebra of subsets of $E$) if
	\begin{enumerate}[(i)]
		\item $E \in \mathcal{E}_0$,
		\item $F \in \mathcal{E}_0 \implies F^c \coloneqq E \setminus F \in \mathcal{E}_0$,
		\item $F, G \in \mathcal{E}_0 \implies F \cup G \in \mathcal{E}_0$.
	\end{enumerate}
\end{defn}

\begin{xtip}{Algebra Properties}{AlgebraProperties}
	Note that $\emptyset = E^c \in \mathcal{E}_0$ and
	\[
		F, G \in \mathcal{E}_0 \implies F \cap G = (F^c \cup G^c)^c \in \mathcal{E}_0.
	\]
\end{xtip}

Thus, an algebra on \(E\) is a family of subsets of \(E\) stable under finitely many set operations.

The collection of all subsets to which we wish to assign a measure is usually a \(\sigma\)-algebra:

\begin{defn}{$\sigma$-Algebra}{}
	A \textit{$\sigma$-algebra} $\mathcal{E}$ on $E$ is a collection of subsets of $E$ that contains $E$ itself, and that is closed under \textit{complements} and \textit{countable unions}; that is:
	\begin{enumerate}
		\item $E \in \mathcal{E}$.
		\item If $A \in \mathcal{E}$, then also $A^c \in \mathcal{E}$.
		\item If $A_1, A_2, \dots \in \mathcal{E}$, then also $\cup_n A_n \in \mathcal{E}$.
	\end{enumerate}
\end{defn}

Thus, a \(\sigma\)-algebra on \(E\) is a family of subsets of \(E\) 'stable under any countable collection of set operations'.

\begin{defn}{Measurable Space}{MeasurableSpace}
	A pair \((E, \cE)\), where \(E\) is a set and \(\cE\) is a \(\sigma\)-algebra on \(E\), is called a \textit{measurable space}. An element of \(\cE\) is called a \(\cE\)-measurable subset of \(E\).
\end{defn}

\begin{defn}{p-System}{}
	A collection $\mathcal{C}$ of subsets of $E$ is called a \textit{p-system} (or $\pi$-system) if it is closed under finite intersections.
\end{defn}

\begin{defn}{d-System}{}
	A collection $\mathcal{D}$ of subsets of $E$ is called a \textit{d-system} (or $\lambda$-system) if it satisfies the following three conditions:
	\begin{enumerate}
		\item $E \in \mathcal{D}$.
		\item If $A, B \in \mathcal{D}$, with $A \supseteq B$, then $A \setminus B \in \mathcal{D}$.
		\item If $(A_n) \in \mathcal{D}$, with $A_1 \subseteq A_2 \subseteq A_3 \subseteq \dots \subseteq \cup_n A_n \coloneqq A$, then $A \in \mathcal{D}$.
	\end{enumerate}
\end{defn}

\begin{thm}{$\sigma$-Algebra as a d- and p-System}{}
	A collection of subsets of $E$ is a $\sigma$-algebra if and only if it is both a p-system and a d-system on $E$.
\end{thm}

\begin{thm}{Monotone Class Theorem}{}
	If a d-system contains a p-system, then it also contains the $\sigma$-algebra generated by that p-system.
\end{thm}

\begin{defn}{Product Space}{}
	The \textit{product space} of two measurable spaces $(E, \cE)$ and $(F, \cF)$ is the measurable space $(E \times F, \cE \otimes \cF)$, where $E \times F \coloneqq \{(x, y) : x \in E, y \in F\}$ is the Cartesian product of $E$ and $F$, and $\cE \otimes \cF$ is the \textit{product $\sigma$-algebra} on $E \times F$ that is generated by the rectangle sets
	\begin{equation*}
		A \times B, \quad A \in \cE, B \in \cF.
	\end{equation*}
\end{defn}

\subsection{Measurable Functions}

Now we have established the basics of measurable spaces, we are going to define a class of "well-behaved" functions between such spaces.

\begin{defn}{Measurable Function}{}
	Let $(E, \cE)$ and $(F, \cF)$ be two measurable spaces, and let $f$ be a function from $E$ to $F$. The function $f$ is called \textit{$\cE/\cF$-measurable} if for all $B \in \cF$ it holds that
	\begin{equation}
		\{x \in E : f(x) \in B\} \in \cE.
		\label{eq:MeasurabilityCondition}
	\end{equation}
\end{defn}

\begin{figure}[H]
	\centering
	\begin{tikzpicture}[
		thick,
		>={Stealth},
		space/.style = {fill=airforceblue!20, draw=black},
		subset/.style = {fill=airforceblue, draw=black},
		label/.style = {font=\small}
		]
		\filldraw[space] (0,0) ellipse (1.6cm and 2.0cm);
		\node[label, above] at (0, 2.1) {$(E, \cE)$};
		\filldraw[subset] plot[smooth cycle, tension=0.7] coordinates {
				(-0.8, -0.2)
				(-0.4, -0.8)
				(0.6, -0.4)
				(0.8, 0.8)
				(0.2, 1.2)
				(-0.6, 0.5)
			};
		\node at (0, 0) {$f^{-1}B$};
		\begin{scope}[shift={(6,0)}]
			\filldraw[space] (0,0) ellipse (2.0cm and 1.2cm);
			\node[label, above] at (0, 1.3) {$(F, \cF)$};
			\filldraw[subset] plot[smooth cycle, tension=0.7] coordinates {
					(-0.6, -0.2)
					(0, -0.5)
					(0.6, -0.1)
					(0.3, 0.5)
					(-0.4, 0.4)
				};
			\node at (0, 0) {$B$};
		\end{scope}
		\draw[->] (0.2, 1.2) to[out=45, in=135] node[midway, above] {$f$} (5.6, 0.5);
	\end{tikzpicture}
	\caption{All elements in $f^{-1}B$ map into $B$.}
	\label{fig:InverseImageMapping}
\end{figure}

\begin{thm}{Measurable Function on a Generated $\sigma$-Algebra}{MeasurableFunctionOnGeneratedSigmaAlgebra}
	A function $f : E \to F$ is $\cE/\cF$-measurable if and only if \Cref{eq:MeasurabilityCondition} holds for all $B \in \cF_0$, with $\cF_0$ generating $\cF$.
\end{thm}

\begin{thm}{Composition of Measurable Functions}{CompositionofMeasurableFunctions}
	Let $(E, \cE)$, $(F, \cF)$, and $(G, \cG)$ be measurable spaces. If $f : E \to F$ and $g : F \to G$ are respectively $\cE/\cF$- and $\cF/\cG$-measurable, then $g \circ f : E \to G$ is $\cE/\cG$-measurable.
\end{thm}

\begin{proof}
	Take any set $B \in \cG$ and let $C \idef \{g \in B\} = g^{-1}(B)$. By the measurability of $g$, $C \in \cF$. Moreover, since $f$ is $\cE/\cF$-measurable, the set $\{f \in C\}$ belongs to $\cE$. But this set is the \textit{same} as $\{g \circ f \in B\}$, as illustrated in \Cref{fig:CompositionMeasurability}.

	\begin{figure}[H]
		\centering
		\begin{tikzpicture}[
			thick,
			>={Stealth},
			space/.style={fill=airforceblue!20, draw=black, shape=ellipse, minimum width=3cm, minimum height=4.4cm},
			subset/.style={fill=airforceblue, draw=black},
			label/.style={font=\small},
			textnode/.style={font=\footnotesize}
			]
			\node[space] (E) at (0,0) {};
			\node[label, above=2pt] at (E.north) {$(E, \cE)$};
			\path[subset, shift={(E.center)}, xscale=1.25]
			plot[smooth cycle, tension=0.7] coordinates {
					(-0.6, -0.4) (-0.3, -0.9) (0.5, -0.5) (0.7, 0.6) (0, 1.0) (-0.7, 0.4)
				};
			\node[textnode] at (E.center) {$(g \circ f)^{-1}B$};

			\node[space] (F) at (5,0) {};
			\node[label, above=2pt] at (F.north) {$(F, \cF)$};
			\path[subset, shift={(F.center)}]
			plot[smooth cycle, tension=0.7] coordinates {
					(-0.5, -0.3) (-0.2, -0.8) (0.6, -0.4) (0.4, 0.7) (-0.5, 0.5)
				};
			\node[textnode] at (F.center) {$g^{-1}B$};

			\node[space] (G) at (10,0) {};
			\node[label, above=2pt] at (G.north) {$(G, \cG)$};
			\path[subset, shift={(G.center)}, scale=0.85]
			plot[smooth cycle, tension=0.7] coordinates {
					(-0.6, -0.4) (0.2, -0.6) (0.5, 0.2) (0, 0.6) (-0.5, 0.3)
				};
			\node[textnode] at (G.center) {$B$};

			\draw[->] (E.east) to[out=15, in=165] node[midway, above] {$f$} (F.west);
			\draw[->] (F.east) to[out=15, in=165] node[midway, above] {$g$} (G.west);
			\draw[->, dashed] (E.north) to[out=30, in=150] node[midway, above] {$g \circ f$} (G.north);
		\end{tikzpicture}
		\caption{Composition of measurability}
		\label{fig:CompositionMeasurability}
	\end{figure}

	Thus, the latter set is a member of $\cE$ for any $B \in \cG$. In other words, $g \circ f$ is $\cE/\cG$-measurable.
\end{proof}

Let $(E, \cE)$ be a measurable space. We will often be interested, especially when defining integrals, in real-valued functions that are $\cE/\cB$-measurable, where $\cB$ is the Borel $\sigma$-algebra in $\R$. In fact, it will be convenient to consider the \textit{extended real line} $\Rbar \idef [-\infty, +\infty]$, adding two ``infinity'' points to the set $\R$, and extending the arithmetic in a natural way, but leaving operations such as $\infty - \infty$ and $\infty/\infty$ undefined. The corresponding Borel $\sigma$-algebra, denoted $\xbar{\cB}$, is generated by the p-system consisting of intervals $[-\infty, r], r \in \R$.

\begin{defn}{Numerical Function}{NumericalFunction}
	Let $(E, \cE)$ be a measurable space.
	\begin{itemize}
		\item A function $f : E \to \R$ is called a \textit{real-valued function} on $E$.
		\item A function $f : E \to \Rbar \coloneqq [-\infty, +\infty]$ is called a \textit{numerical function} on $E$.
		\item An $\cE/\xbar{\cB}$-measurable numerical function $f$ is said to be \textit{$\cE$-measurable}; we write $f \in \cE$. Also, $\cE_+$ is the class of \textit{positive $\cE$-measurable functions}.
	\end{itemize}
\end{defn}

We can further expand the collection of measurable functions that can be obtained from simple functions by taking limits of sequences of measurable functions. Let us first recall some definitions regarding sequences of real numbers.

\begin{defn}{Infimum and Supremum}{InfimumandSupremum}
	Let $(a_n)$ be a sequence of real numbers.
	\begin{enumerate}
		\item The \textit{infimum}, $\inf a_n$, is the greatest element in $\Rbar$ that is less than or equal to all $a_n$. It can be $-\infty$.
		\item The \textit{supremum}, $\sup a_n$, is the smallest element in $\Rbar$ that is greater than or equal to all $a_n$. It can be $\infty$.
	\end{enumerate}
\end{defn}

If for some $N \in \N$ it holds that $\inf a_n = a_N$, then we can write $\inf a_n$ also as $\min a_n$. In that case, the set of indexes for which the minimum is attained is denoted by $\argmin a_n$. In the same way, when the supremum is attained by some $N \in \argmax a_n$, we can write $\max a_n$ instead of $\sup a_n$. Similarly, the infimum and supremum can be defined for any set $A \subseteq \Rbar$. In particular, if $A \idef \{a_n, n \in \N\}$, then $\inf A = \inf a_n$.

\begin{defn}{Liminf and Limsup}{}
	Let $(a_n)$ be a sequence of real numbers.
	\begin{enumerate}
		\item The \textit{limit inferior}, $\liminf a_n$, is the eventual lower bound for the sequence $(a_n)$; that is, $\liminf a_n \coloneqq \lim_{m \to \infty} \inf_{n \ge m} a_n = \sup_m \inf_{n \ge m} a_n$.
		\item The \textit{limit superior}, $\limsup a_n$, is the eventual upper bound for the sequence $(a_n)$; that is, $\limsup a_n \coloneqq \lim_{m \to \infty} \sup_{n \ge m} a_n = \inf_m \sup_{n \ge m} a_n$.
	\end{enumerate}
\end{defn}

For a sequence $(f_n)$ of numerical functions on a set $E$, we can define functions
\begin{equation*}
	\inf f_n, \quad \sup f_n, \quad \liminf f_n, \quad \limsup f_n
\end{equation*}
\textit{pointwise}, considering for each $x$ the sequence $(a_n)$ defined by $a_n \idef f_n(x)$.

Obviously, $\liminf f_n \le \limsup f_n$ (pointwise). If the two are equal, we write $\lim f_n$ for the limit, and write $f_n \to f$. If $(f_n)$ is a sequence of functions that increases pointwise, then the limit $f \idef \lim f_n$ always exists. In this case we write $f_n \uparrow f$. A similar notation $f_n \downarrow f$ holds for a decreasing sequence of functions.

\begin{prop}{Limits of Measurable Functions}{}
	Let $(f_n)$ be a sequence of $\cE$-measurable numerical functions. Then, the numerical functions $\sup f_n, \inf f_n, \limsup f_n$, and $\liminf f_n$ are also $\cE$-measurable. In particular, limits of measurable functions (when $\limsup f_n = \liminf f_n$) are measurable as well.
\end{prop}

\begin{proof}
	\begin{enumerate}
		\item Let $f \idef \sup f_n$. Since sets of the form $[-\infty, r]$ form a p-system that generates the Borel $\sigma$-algebra on $\Rbar$, it suffices to show, by \Cref{thm:MeasurableFunctiononaGeneratedSigmaAlgebra}, that $f^{-1}[-\infty, r] \in \cE$ for all $r$. The latter follows from the observation that
		      \begin{equation*}
			      f^{-1}[-\infty, r] = \cap_n f_n^{-1}[-\infty, r],
		      \end{equation*}
		      where $f_n^{-1}[-\infty, r] \in \cE$ by the measurability of $f_n$, and the fact that $\cE$ is closed under countable intersections.
		\item To show $\inf f_n \in \cE$, we simply observe that $\inf f_n = - \sup(-f_n)$.
		\item Since $\limsup f_n = \inf_m \sup_{n \ge m} f_n$, we can use Steps 1 and 2 above to conclude that $\limsup f_n \in \cE$.
		\item Similarly, since $\liminf f_n = \sup_m \inf_{n \ge m} f_n$, we have $\liminf f_n \in \cE$.
	\end{enumerate}
\end{proof}

\begin{thm}{Approximation from Below}{}
	A \textit{positive} numerical function on $E$ is $\cE$-measurable if and only if it is the (pointwise) limit of an increasing sequence of simple functions.
\end{thm}

\begin{thm}{Positive and Negative Part of a Function}{}
	A function $f$ is $\cE$-measurable if and only if both $f^+$ and $f^-$ are.
\end{thm}

\begin{proof}
	To prove sufficiency, suppose $f^+$ and $f^-$ are measurable. Consider the decomposition illustrated in \Cref{fig:PositiveNegativeParts}.

	\begin{figure}[H]
		\centering
		\begin{tikzpicture}[
			thick,
			>={Stealth},
			axis/.style={->, draw=black!60, thin},
			func/.style={draw=airforceblue, line width=1pt},
			guide/.style={draw=black!30, dashed, thin},
			label/.style={font=\footnotesize, anchor=east}
			]

			\begin{scope}[yshift=0cm]
				\draw[axis] (-3,0) -- (3,0) node[right] {$x$};
				\draw[axis] (0,-1.2) -- (0,1.5) node[above] {$f(x)$};
				\draw[func] plot[smooth, tension=0.6] coordinates {
						(-2.5, -0.4) (-1.8, 0) (-0.6, 1.2) (0.6, 0) (1.6, -0.8) (2.4, 0) (2.8, 0.2)
					};
				\node at (2, 1.2) {$f$};
			\end{scope}

			\begin{scope}[yshift=-3.5cm]
				\draw[axis] (-3,0) -- (3,0) node[right] {$x$};
				\draw[axis] (0,-0.2) -- (0,1.5) node[above] {$f^+(x)$};
				\draw[func] (-2.5, 0) -- (-1.8, 0);
				\draw[func] (0.6, 0) -- (2.4, 0);
				\draw[func] plot[smooth, tension=0.6] coordinates { (-1.8, 0) (-0.6, 1.2) (0.6, 0) };
				\draw[func] plot[smooth, tension=0.6] coordinates { (2.4, 0) (2.8, 0.2) };
				\draw[guide] plot[smooth, tension=0.6] coordinates { (-2.5, -0.4) (-1.8, 0) };
				\draw[guide] plot[smooth, tension=0.6] coordinates { (0.6, 0) (1.6, -0.8) (2.4, 0) };
				\node at (2, 1.2) {$f^+ \idef \max(f, 0)$};
			\end{scope}

			\begin{scope}[yshift=-7cm]
				\draw[axis] (-3,0) -- (3,0) node[right] {$x$};
				\draw[axis] (0,-0.2) -- (0,1.5) node[above] {$f^-(x)$};
				\draw[func] (-1.8, 0) -- (0.6, 0);
				\draw[func] (2.4, 0) -- (2.8, 0);
				\draw[func] plot[smooth, tension=0.6] coordinates { (-2.5, 0.4) (-1.8, 0) };
				\draw[func] plot[smooth, tension=0.6] coordinates { (0.6, 0) (1.6, 0.8) (2.4, 0) };
				\draw[guide] plot[smooth, tension=0.6] coordinates { (-1.8, 0) (-0.6, 1.2) (0.6, 0) };
				\draw[guide] plot[smooth, tension=0.6] coordinates { (2.4, 0) (2.8, 0.2) };
				\node at (2, 1.2) {$f^- \idef -\min(f, 0)$};
			\end{scope}
		\end{tikzpicture}
		\caption{Decomposition of a function into positive and negative parts. Both $f^+$ and $f^-$ are non-negative.}
		\label{fig:PositiveNegativeParts}
	\end{figure}

	We observe the following: For $r \ge 0$, we have $\{f \le r\} = \{f^+ \le r\}$. For $r < 0$, we have $\{f \le r\} = \{f^- \ge -r\}$. In both cases the inverse image of $f$ is in $\cE$. Hence, $f$ is $\cE$-measurable.

	Necessity is proved in a similar way; e.g., $\{f^+ \le r\} = \{f \le r\} \in \cE$ for $r \ge 0$ and is equal to the empty set (also in $\cE$) for $r < 0$.
\end{proof}

The important point of the preceding results is that every measurable function can be viewed as the pointwise limit of simple functions, which themselves are constructed as linear combinations of indicator functions. As a consequence, arithmetic operations on $\cE$-measurable functions $f$ and $g$, such as $f + g$, $f - g$, $fg$, $f/g$, $f \wedge g$, and $f \vee g$ are $\cE$-measurable as well, as long as these functions are well-defined (for example, $\infty - \infty$ and $\infty/\infty$ are undefined).

As an example, if $f, g \in \cE$, then $f^+, f^-, g^+$ and $g^-$ are in $\cE_+$. Moreover, each of these functions is the limit of a sequence of simple functions $f_n^+, f_n^-, g_n^+$, and $g_n^-$, respectively. Thus, $f - g$ is the limit of the simple function $f_n^+ - f_n^- - g_n^+ + g_n^-$, which lies in $\cE$. Hence, the limit is also in $\cE$.

We conclude this section with the notion of a monotone class of functions.

\begin{defn}{Monotone Class of Functions}{}
	Let $(E, \cE)$ be a measurable space. A collection $\cM$ of numerical functions on $E$ is called a \textit{monotone class} if it satisfies:
	\begin{enumerate}
		\item $\I_E \in \cM$.
		\item If $f, g \in \cM$ are bounded and $a, b \in \R$, then $af + bg \in \cM$.
		\item If $(f_n)$ is a sequence of positive functions in $\cM$ that increases to some $f$, then $f \in \cM$.
	\end{enumerate}
\end{defn}

\begin{thm}{Monotone Class Theorem for Functions}{MonotoneClassTheoremforFunctions}
	Let $\cM$ be a monotone class of functions on $E$ and let $\mathcal{C}$ be a p-system that generates $\cE$. If $\I_A \in \cM$ for every $A \in \mathcal{C}$, then $\cM$ includes
	\begin{itemize}
		\item all positive $\cE$-measurable functions,
		\item all bounded $\cE$-measurable functions.
	\end{itemize}
\end{thm}

\subsection{Measures}
Let \((E,\cE)\) be a measurable space. We wish to assign a measure to all sets in \(\cE\) that has the same properties as an area or length measure. Particularly, any measure should have the property that the measure of the union of a disjoint (i.e., non-overlapping) collection of sets is equal to the sum of the measures for the individual sets, as illustrated in \Cref{fig:MeasureProperties}.

\begin{figure}[H]
	\centering
	\begin{tikzpicture}[
			thick,
			line join=round,
			tri/.style={draw=black, fill=airforceblue!50}
		]
		\def\s{1}
		\begin{scope}[scale=\s]
			\def\tris{{(-2.6,1.4)--(0.6,3.2)--(0.9,0.8)},{(1.1,2.2)--(2.2,3.0)--(2.2,1.2)},{(1.1,1.0)--(1.9,0.8)--(1.5,0.0)},{(-1.8,-0.6)--(-0.6,0.6)--(-0.6,-1.4)},{(-0.3,0.3)--(3.0,-0.2)--(0.8,-2.0)}}
			\foreach \t in \tris { \path[tri] \t -- cycle; }
		\end{scope}
	\end{tikzpicture}
	\caption{Any measure has the same properties as the area measure. For example, the total area of the non-overlapping triangles is the sum of the areas of the individual triangles.}
	\label{fig:MeasureProperties}
\end{figure}

\begin{defn}{Measure}{Measure}
	Let $(E, \cE)$ be a measurable space. A \textit{measure} $\mu$ on $\cE$ is a set function
	\begin{equation*}
		\mu : \cE \to [0, \infty]
	\end{equation*}
	with $\mu(\emptyset) = 0$, such that for every sequence $(A_n)$ of disjoint sets in $\cE$,
	\begin{equation}
		\mu(\cup_n A_n) = \sum_n \mu(A_n).
		\label{eq:SigmaAdditivity}
	\end{equation}
\end{defn}


\begin{thm}{Carathéodory's Extension Theorem}{CarathéodorysExtensionTheorem}
	Let $E$ be a set with an algebra $\cE_0$ thereon. Every pre-measure $\mu_0 : \cE_0 \to [0, \infty]$ can be extended to a measure $\mu$ on $(E, \sigma(\cE_0))$.
\end{thm}

\begin{lem}{Elementary Inequalities}{ElementaryInequalities}
	Let $(E, \cE, \mu)$ be a measure space. Then
	\begin{enumerate}[label=(\alph*)]
		\item $\mu(A \cup B) \le \mu(A) + \mu(B) \quad (A, B \in \cE)$,
		\item $\mu\left(\bigcup_{i \le n} F_i\right) \le \sum_{i \le n} \mu(F_i) \quad (F_1, F_2, \dots, F_n \in \cE)$.
	\end{enumerate}
	Furthermore, if $\mu(E) < \infty$, then
	\begin{enumerate}[label=(\alph*), resume]
		\item $\mu(A \cup B) = \mu(A) + \mu(B) - \mu(A \cap B) \quad (A, B \in \cE)$,
		\item (inclusion-exclusion formula): for $F_1, F_2, \dots, F_n \in \cE$,
		      \begin{align*}
			      \mu\left(\bigcup_{i \le n} F_i\right) & = \sum_{i \le n} \mu(F_i) - \sum \sum_{i < j \le n} \mu(F_i \cap F_j)                                                     \\
			                                            & + \sum \sum \sum_{i < j < k \le n} \mu(F_i \cap F_j \cap F_k) - \dots + (-1)^{n-1} \mu(F_1 \cap F_2 \cap \dots \cap F_n),
		      \end{align*}
		      successive partial sums alternating between over- and under-estimates.
	\end{enumerate}
\end{lem}

\begin{prop}{Properties of a Measure}{}
	Let $\mu$ be a measure on a measurable space $(E, \cE)$, and let $A, B$ and $A_1, A_2, \dots$ be measurable sets. Then, the following hold:
	\begin{enumerate}
		\item (\textit{Monotonicity}): $A \subseteq B \implies \mu(A) \le \mu(B)$.
		\item (\textit{Continuity from below}): If $A_1, A_2, \dots$ is increasing (i.e., $A_1 \subseteq A_2 \subseteq \cdots$), then
		      \begin{equation*}
			      \lim_n \mu(A_n) = \mu(\cup_n A_n).
		      \end{equation*}
		\item (\textit{Countable subadditivity}): $\mu(\cup_n A_n) \le \sum_n \mu(A_n)$.
	\end{enumerate}
\end{prop}

\begin{thm}{Uniqueness of Finite Measures}{UniquenessofFiniteMeasures}
	Let $\mathcal{C}$ be a p-system that generates $\cE$. If $\mu$ and $\nu$ are two measures on $(E, \cE)$ with $\mu(E) = \nu(E) < \infty$, then
	\begin{equation*}
		\mu(A) = \nu(A) \text{ for all } A \in \mathcal{C} \implies \mu(A) = \nu(A) \text{ for all } A \in \cE.
	\end{equation*}
\end{thm}

\begin{defn}{Product Measure}{}
	Let $(E, \cE, \mu)$ and $(F, \cF, \nu)$ be two measure spaces. The \textit{product measure} of $\mu$ and $\nu$ on the product space $(E \times F, \cE \otimes \cF)$ is defined as the measure $\pi$ with
	\begin{equation}
		\pi(A \times B) \coloneqq \mu(A) \nu(B), \quad A \in \cE, B \in \cF.
		\label{eq:ProductMeasureDef}
	\end{equation}
	This measure $\pi$ is often written as $\mu \otimes \nu$.
\end{defn}

\subsection{Integrals}
We combine measurable functions and measures to define integrals in a unified and elegant way. We combine measurable functions and measures to define integrals
in a unified and elegant way. The procedure for building up proofs for measurable functions from indicator functions, simple functions, and limits thereof is often referred to as “the standard machine”.

\begin{prop}{Insensitivity of the Integral}{}
	Let $(E, \cE, \mu)$ be a measure space and let $f$ and $g$ be positive $\cE$-measurable functions that are equal $\mu$-almost everywhere. Then, $\mu f = \mu g$.
\end{prop}

\begin{thm}{Properties of an Integral}{}
	Below $a, b$ are in $\R_+$ and $f, g, f_n$ are in $\cE_+$.
	\begin{enumerate}
		\item (\textit{Positivity and Monotonicity}): $\mu f \ge 0$, and $f = 0 \implies \mu f = 0$. Also, $f \le g \implies \mu f \le \mu g$.
		\item (\textit{Linearity}): $\mu(af + bg) = a \mu f + b \mu g$.
		\item (\textit{Monotone convergence}): If $f_n \uparrow f$, then $\mu f_n \uparrow \mu f$.
	\end{enumerate}
\end{thm}

\begin{thm}{Measures from Linear Functionals}{}
	Let $L : \cE_+ \to \Rbar_+$. Then, there exists a unique measure $\mu$ on $(E, \cE)$ such that $L(f) = \mu f$ if and only if
	\begin{enumerate}
		\item $f = 0 \implies L(f) = 0$,
		\item $f, g \in \cE_+$ and $a, b \in \R_+ \implies L(af + bg) = aL(f) + bL(g)$,
		\item $(f_n) \in \cE_+$ and $f_n \uparrow f \implies L(f_n) \uparrow L(f)$.
	\end{enumerate}
\end{thm}

\begin{thm}{Integrals with Densities}{}
	Let $\nu \coloneqq p \mu$. Then, for every $f \in \cE_+$, we have $\nu f = \mu(pf)$.
\end{thm}

\begin{thm}{Radon--Nikodym}{}
	Let $\nu$ and $\mu$ be two measures on $(E, \cE)$ such that $\mu$ is $\sigma$-finite and $\nu \ll \mu$. Then, there exists a $p \in \cE_+$ such that $\nu = p \mu$. Moreover, $p$ is unique up to equivalence; that is, if there is another $\tilde{p} \in \cE_+$ such that $\nu = \tilde{p} \mu$, then $p = \tilde{p}$, $\mu$-almost everywhere.
\end{thm}

\begin{thm}{Integration with Respect to Image Measures}{}
	Let $\nu = \mu \circ h^{-1}$ be the image measure of $\mu$ under $h$. Then, for every positive $\cF$-measurable function $f$,
	\begin{equation*}
		\nu f = (\mu \circ h^{-1}) f = \int_F f \di \nu = \int_E f \circ h \di \mu = \mu(f \circ h).
	\end{equation*}
\end{thm}

\begin{defn}{Transition Kernel}{}
	Let $(E, \cE)$ and $(F, \cF)$ be measurable spaces. A \textit{transition kernel} from $(E, \cE)$ into $(F, \cF)$ is a mapping $K : E \times \cF \to \Rbar_+$ such that:
	\begin{enumerate}
		\item $K(\cdot, B)$ is $\cE$-measurable for every $B \in \cF$.
		\item $K(x, \cdot)$ is a measure on $(F, \cF)$ for every $x \in E$.
	\end{enumerate}
	If $(F, \cF) = (E, \cE)$ and $K(x, \cdot)$ is a probability measure for every $x \in E$, $K$ is said to be a \textit{probability transition kernel} on $(E, \cE)$.
\end{defn}

\begin{thm}{Measures and Functions from Kernels}{}
	Let $K$ be a transition kernel from $(E, \cE)$ into $(F, \cF)$.
	\begin{enumerate}
		\item For all $f \in \cF_+$ the function $Kf$ defined by $(Kf)(x) \coloneqq \int_F K(x, \di y) f(y)$, $x \in E$, is in $\cE_+$.
		\item For every measure $\mu$ on $(E, \cE)$ the set function $\mu K$ defined by $(\mu K)(B) \coloneqq \int_E \mu(\di x) K(x, B)$, $B \in \cF$, is a measure on $(F, \cF)$.
		\item For every measure $\mu$ on $(E, \cE)$ and $f \in \cF_+$,
		      \begin{equation*}
			      (\mu K) f = \mu (K f) = \int_E \mu(\di x) \int_F K(x, \di y) f(y).
		      \end{equation*}
	\end{enumerate}
\end{thm}

\begin{thm}{Fubini}{}
	Let $\mu$ and $\nu$ be $\Sigma$-finite measures on $(E, \cE)$ and $(F, \cF)$, respectively.
	\begin{enumerate}
		\item There exists a unique $\Sigma$-finite measure $\pi$ on the product space $(E \times F, \cE \otimes \cF)$ such that for every positive $f$ in $\cE \otimes \cF$,
		      \begin{equation}
			      \pi f = \int_E \mu(\di x) \int_F \nu(\di y) f(x, y) = \int_F \nu(\di y) \int_E \mu(\di x) f(x, y).
			      \label{eq:FubiniIntegral}
		      \end{equation}
		\item If $f \in \cE \otimes \cF$ and is $\pi$-integrable, then $f(x, \cdot)$ is $\nu$-integrable for $\mu$-almost every $x$, and $f(\cdot, y)$ is $\mu$-integrable for $\nu$-almost every $y$ and \Cref{eq:FubiniIntegral} holds again.
	\end{enumerate}
\end{thm}
